% !TEX root = ../main_ver2_bench_only.tex
\section{Introduction}
\label{sec:intro}

Large language model (LLM) agents increasingly rely on external \emph{skills}---structured tool descriptions that extend model capabilities beyond parametric knowledge---to perform domain-specific tasks.
Recent work demonstrates that curated skill libraries substantially improve agent performance: \citet{skillsbench2026} report an average gain of 16.2 percentage points across domains, with the medical domain exhibiting the largest improvement at 51.9 percentage points.
This finding underscores the critical role of skill quality in high-stakes settings such as clinical decision support, biomedical literature analysis, and drug safety assessment.

Despite this promise, medical skills present unique evaluation challenges that existing benchmarks do not address.
Medical skill specifications encode domain-specific safety boundaries, clinical workflow constraints, and a fundamental \emph{reasoning--execution duality}: some skills wrap executable tools (BLAST, RDKit, clinical calculators), while others provide purely reasoning-based guidance (differential diagnosis protocols, evidence synthesis frameworks).
Both modes are clinically legitimate but require fundamentally different quality criteria.

\paragraph{Task-level evaluation obscures skill-level failures.}
Existing medical agent benchmarks evaluate end-to-end task success \citep{healthbench2025, medagentbench2025, labbench2024} but do not isolate whether the agent correctly \emph{executed} the skill specification it was given.
An agent may answer a clinical question correctly by ignoring a poorly written drug interaction checker, or it may fail a genomic analysis task despite faithfully following a specification that omits critical parameter constraints.
Task-level metrics conflate these fundamentally different failure modes, leaving skill developers without actionable feedback on \emph{where} in the specification--execution chain failures originate.

\medskillbench addresses this gap through systematic skill-level evaluation that treats skill descriptions as first-class agent interfaces whose execution quality can be independently measured, decomposed across eight diagnostic dimensions, and stratified by medical domain.\footnote{Code, data, and evaluation artifacts are available at \anonrepo.}

\begin{figure}[t]
\figbox{5.2cm}{Figure 1 (Teaser): the \rbu gap. Left panel---agents read \skillmd at 96.2\% invocation. Right panel---only 74.3\% complete the specified task. The 22-pp wedge between the two represents specification grounding failure.}
\caption{The \rbu gap: agents read skill specifications (96.2\%) yet only complete skill-grounded tasks at 74.3\%, exposing a 22-percentage-point specification grounding failure.}
\label{fig:teaser}
\end{figure}

\paragraph{Contribution 1: \medskillbench.}
We introduce the first large-scale skill-level execution benchmark for medical AI agents.
\medskillbench aggregates 1,462 skills from three open-source medical skill repositories and evaluates each in an isolated sandbox environment.
A mandatory \emph{read-first} protocol ensures the agent must parse the \skillmd specification before attempting execution.
An LLM-as-judge scores each execution trajectory along eight binary dimensions---\emph{skill invocation}, \emph{correct usage}, \emph{task completion}, \emph{tool execution success}, \emph{output quality}, \emph{reasoning quality}, \emph{efficiency}, and \emph{trajectory quality}---providing fine-grained diagnostic signal rather than a single pass/fail label.

\paragraph{Contribution 2: \rbu gap quantification.}
Using \medskillbench, we identify and quantify a systematic \emph{read-but-not-use} (\rbu) gap: agents invoke skills at a rate of 96.2\% but complete the specified tasks at only 74.3\%, yielding a 22-percentage-point gap.
This gap represents a \emph{specification grounding failure}---the model recognizes what a skill describes but cannot reliably translate that description into correct execution.
Stratifying by \texttt{tool\_type} reveals that the gap is strongly modulated by skill interface design: guide-style skills (no executable tool) exhibit an \rbu gap of 0.269, while tool-based skills show substantially smaller gaps (CLI: 0.106; Python: 0.099; mixed: 0.069; R: 0.111).
This finding shifts the diagnosis from ``models cannot use tools'' to ``models struggle most when skill descriptions lack executable grounding.''
A formal two-stage decomposition reveals that the gap is overwhelmingly driven by semantic grounding failure ($\delta_{\text{sem}} = 0.242$) rather than execution failure ($\delta_{\text{exec}} = -0.023$), with the negative execution component indicating that models sometimes succeed by bypassing the skill specification entirely.

\paragraph{Contribution 3: Skill-augmented QA evaluation.}
We evaluate whether skill availability measurably improves model performance on real medical benchmark questions.
By pairing validated skill--question matches with a controlled with/without-skill ablation across multiple executor models, we quantify the accuracy uplift attributable to skill specifications and show that the uplift is strongly modulated by skill interface type---consistent with the \rbu gap analysis.
\todo{Report uplift magnitude and statistical significance.}

Our work is organized around three research questions:
\begin{enumerate}[nosep,leftmargin=*,label=\textbf{RQ\arabic*:}]
    \item Do medical agent skills exhibit a systematic read-but-not-use gap between specification comprehension and execution?
    \item Is the gap primarily attributable to environment dependency or to description-level grounding failure?
    \item Does skill availability improve model accuracy on real medical QA benchmarks, and is the uplift modulated by skill execution quality?
\end{enumerate}

\paragraph{Positioning.}
Following the taxonomy of \citet{skillsbench2026}, \medskillbench operates primarily at the \emph{Skill Deployment} tier (T2), evaluating skills post-authoring.
It secondarily informs \emph{Skill Development} (T1) by identifying which description patterns yield higher groundability.
Unlike end-to-end medical benchmarks and biomedical agent systems \citep{medagentbench2025,labbench2024,reflectool2025,biomni2025,stella2025}, it does not propose another general medical agent scaffold; it targets the interface layer between authored skill specifications and grounded execution.
The benchmark requires no model training and is applicable to closed-weight models and resource-constrained settings alike.

The remainder of this paper is structured as follows.
\S\ref{sec:related_work} reviews related work.
\S\ref{sec:method} details the \medskillbench evaluation protocol and the skill-augmented QA evaluation design.
\S\ref{sec:experiments} reports experimental results, including skill-augmented performance analysis, and discusses implications.
\S\ref{sec:conclusion} concludes.
